\scanpage{96}
% archive foliation: 95
\uline{1.5.}\quad\uline{Sections horizontales de} $P$

Pour quelles connexions projectives sur $X/S$ (avec $\underline{\omega}_{X/S}$ inv.) y a-t-il des sections horizontales \inserted{de $P$,} et quelles sont-elles ? On va étudier d'abord les sections horizontales de $P$ là où elles sont disjointes de la section donnée $\sigma$. Elles correspondent donc à une décomp. de $\mathcal{L}$ en somme directe, ou de $\mathcal{H}$ en
\[
\mathcal{H}\simeq\Omega\oplus\Omega^0\oplus\Omega^{-1}\qquad(\Omega\simeq\omega)
\]
\deleted{Si la \unsure{connexion}} Une connexion $c$ \unsure{donne} la section $\sigma'$ horizontale \unsure{ssi} \deleted{$c(\Omega^{-1})\subset$}
\[
c(\Omega^{-1})\subset(\Omega^0+\Omega^{-1})\otimes\underline{\omega}
\]
\inserted{ou $c(\Omega^0)\subset(\Omega^0+\Omega^{-1})\otimes\underline{\omega}$}
\margin{$\uline{2\text{ inv.!}}$}
\unsure{comme la donnée} d'une \deleted{section}\inserted{base} $\delta$ de $\Omega$, permettra de décrire $c$ par trois formes $\bar\omega_2$, $\bar\omega_1$, $\bar\omega_0$, à valeurs dans $\Omega\otimes\underline{\omega}$, $\Omega^0\otimes\underline{\omega}=\underline{\omega}$ et $\Omega^{-1}\otimes\underline{\omega}$ :
\[
-\unsure{2}\bar\omega_2=0
\]
\deleted{donc si $2$ inversible ceci signifie $\bar\omega_2=0$} et peut donc s'obtenir (si un \unsure{choix} de connexion \unsure{mod} $\Gamma(\Omega\otimes\omega)$ est donné) \deleted{\unsure{précis.}} \inserted{\unsure{préc.}} \unsure{par la} restriction à $\Omega$ -- puis on obtient de façon unique en corrigeant $c$ par adjonction d'une $\omega_2\in\Gamma\Omega\otimes\underline{\omega}$.

Donc

\uline{Prop.} Pour tout splitting de \inserted{l'ext.} $\mathcal{L}$ ($\simeq P^1_{X/S}\underline{\omega}\otimes\underline{\omega}^{-1}$ dans le cas ``canonique'') il y a une et une seule \uline{conn. admissible} sur $\mathcal{H}$ (ou $P$ --) pour laquelle la section correspondante de $P$ soit horizontale. Donc \deleted{\unsure{caractérisation bien}} \deleted{une} \deleted{injection}\inserted{application} de l'ens. des sections de $P-\sigma(X)$ sur $X$, \inserted{i.e. des splittings de $P^1_{X/S}(\omega)$,} dans l'ens. des connexions \deleted{\unsure{différentielles}} sur $X$, ayant comme image celles qui admettent une section hor. disjointe de $\sigma$.
